\chapter{Test Fixture Specification}\label{ch:test-fixtures}\label{sim:ch:test-fixtures}

This chapter specifies the test fixture requirements for all domains, addressing
the review finding that ``testing sections lack substance.'' Fixtures provide
standardised test data for unit, integration, and performance testing.

\section{Overview}

Test fixtures are pre-defined datasets that enable:
\begin{itemize}
  \item Deterministic test execution (same input $\rightarrow$ same output)
  \item Coverage of edge cases and boundary conditions
  \item Regression testing across releases
  \item Performance benchmarking with consistent data volumes
\end{itemize}

All fixtures \textbf{must} be sanitised for PII and checked into version control.

\section{Fixture Count Calculation Methodology}
\label{sec:fixture-count-methodology}

The fixture counts specified in this chapter are derived systematically, not
arbitrarily. This section documents the calculation methodology for auditability
and future extension.

\subsection{Coverage-Based Calculation}

Fixture counts are derived using a \textbf{pairwise coverage} approach (2-way
combinatorial testing), which ensures that every pair of input parameters is
tested at least once. For a feature with $k$ parameters, each with $v_i$ values:

\[
\text{Min Fixtures} = \max\left(\prod_{i=1}^{2} v_{\text{top-}i}\right)
\]

\noindent
where $v_{\text{top-}1}$ and $v_{\text{top-}2}$ are the two parameters with the
most values. This provides a lower bound; actual counts include additional
edge case and boundary fixtures.

\subsection{Fixture Count Derivation by Domain}

\begin{table}[htbp]
\centering
\caption{Fixture count derivation methodology.}
\label{tab:fixture-derivation}
\footnotesize
\begin{tabularx}{\textwidth}{@{}l l X@{}}
\toprule
\textbf{Domain} & \textbf{Key Parameters} & \textbf{Derivation} \\
\midrule
Arabic AP Invoices & 
  Country (4), VAT category (5), Invoice type (4) &
  $4 \times 5 = 20$ standard + 35 edge/invalid \\
TB Extracts &
  Entity type (3), Period type (4), Volume (3) &
  $3 \times 4 = 12$ standard + additional volumes \\
Simula Training &
  Complexity (3), SQL type (5), Schema size (3) &
  $5 \times 3 = 15$ per category, 4.6x for coverage \\
\bottomrule
\end{tabularx}
\end{table}

\subsection{Edge Case Multiplier}

Beyond combinatorial coverage, fixtures include:
\begin{description}
  \item[Boundary values (1.5x):] Min, max, and mid-range values for numeric fields
  \item[Invalid inputs (0.5x):] Each valid category has corresponding invalid cases
  \item[Null/empty handling (0.2x):] Optional fields tested with missing values
\end{description}

\noindent
\textbf{Example}: Arabic AP standard fixtures (20) $\times$ 2.2 multiplier $\approx$ 44,
rounded to 55 to ensure coverage gaps are addressed.

\subsection{Extending Fixture Counts}

When new features require additional fixtures:
\begin{enumerate}
  \item Identify new parameters introduced by the feature
  \item Calculate 2-way coverage requirement: $\text{new} = \max(v_{\text{new}} \times v_{\text{existing}})$
  \item Apply edge case multiplier (typically 2.0--2.5x)
  \item Document the derivation in the PR adding the fixtures
\end{enumerate}

\section{Fixture organisation}

\begin{lstlisting}[basicstyle=\ttfamily\small,caption={Test fixture directory structure.}]
tests/
  fixtures/
    arabic/
      invoices/
        sample_invoices.json       # Standard test invoices
        edge_cases.json            # Boundary conditions
        invalid_invoices.json      # Validation failure cases
      ocr/
        ocr_outputs.json           # OCR extraction results
        confidence_levels.json     # Various confidence scores
      vat/
        vat_rules.json             # Rule definitions
        expected_checklists.json   # Expected outputs
    tb/
      extracts/
        sample_tb_extract.json     # Standard TB data
        multi_entity.json          # Multi-legal entity
      variances/
        sample_variances.json      # Normal variances
        anomalies.json             # Anomaly test cases
      commentary/
        expected_commentary.json   # AI output targets
    simula/
      taxonomy/
        sample_taxonomy.json       # Taxonomy tree
        complexity_levels.json     # Complexity examples
      examples/
        training_examples.json     # Sample training data
        rejected_examples.json     # Rejection cases
    common/
      schemas/                     # Schema copies for offline validation
      mock_responses/              # Mock LLM responses
\end{lstlisting}

\section{Arabic AP fixtures}

\subsection{Invoice fixtures}

\begin{table}[htbp]
\centering
\caption{Required invoice fixture categories.}
\label{tab:invoice-fixtures}
\footnotesize
\begin{tabularx}{\linewidth}{@{}l l X@{}}
\toprule
\textbf{Category} & \textbf{Count} & \textbf{Coverage} \\ \midrule
Standard invoices & 20 & All four countries (SA, EG, BH, QA) \\
Edge cases & 15 & Missing fields, unusual formats \\
Invalid invoices & 10 & Schema violations, invalid TINs \\
Credit notes & 5 & Negative amounts, adjustments \\
Multi-currency & 5 & Non-local currency invoices \\
\bottomrule
\end{tabularx}
\end{table}

\subsubsection{Fixture schema}

Each invoice fixture \textbf{must} include:
\begin{lstlisting}[basicstyle=\ttfamily\small]
{
  "fixture_id": "INV-SA-001",
  "description": "Standard Saudi invoice with 15% VAT",
  "input": { /* invoice.schema.json conformant */ },
  "expected_output": {
    "vat_category": "FI",
    "checklist_results": ["Y", "Y", "Y", ...],
    "verdict": "RELEASE"
  },
  "tags": ["saudi", "standard", "vat-15"]
}
\end{lstlisting}

\subsection{OCR fixtures}

\begin{table}[htbp]
\centering
\caption{OCR fixture confidence levels.}
\label{tab:ocr-fixtures}
\footnotesize
\begin{tabularx}{\linewidth}{@{}l r X@{}}
\toprule
\textbf{Category} & \textbf{Confidence} & \textbf{Purpose} \\ \midrule
High quality & $>$85\% & Happy path testing \\
Medium quality & 55--85\% & Threshold boundary testing \\
Low quality & 30--55\% & Manual review trigger \\
Very low quality & $<$30\% & Rejection testing \\
Mixed quality & Varies & Field-level confidence variation \\
\bottomrule
\end{tabularx}
\end{table}

\section{Trial Balance fixtures}

\subsection{TB extract fixtures}

\begin{table}[htbp]
\centering
\caption{Required TB extract fixture categories.}
\label{tab:tb-fixtures}
\footnotesize
\begin{tabularx}{\linewidth}{@{}l l X@{}}
\toprule
\textbf{Category} & \textbf{Count} & \textbf{Coverage} \\ \midrule
Standard TB & 10 & Various account types, normal balances \\
Multi-entity & 5 & Inter-company, consolidation \\
Period variations & 5 & Month-end, quarter-end, year-end \\
Large volume & 3 & 10K, 50K, 100K accounts \\
Anomalous & 10 & Various anomaly types \\
\bottomrule
\end{tabularx}
\end{table}

\subsection{Variance fixtures}

Each variance fixture \textbf{must} cover:
\begin{itemize}
  \item \textbf{Normal variance}: Within materiality threshold
  \item \textbf{Material variance}: Exceeds threshold, no anomaly
  \item \textbf{Spike anomaly}: Period-over-period $>3\sigma$
  \item \textbf{Sign flip}: Balance sign reversal
  \item \textbf{New account}: Account appearing for first time
  \item \textbf{Stale balance}: Unchanged for $>$6 months
\end{itemize}

\section{Simula fixtures}

\subsection{Taxonomy fixtures}

\begin{lstlisting}[basicstyle=\ttfamily\small,caption={Taxonomy fixture structure.}]
{
  "fixture_id": "TAX-001",
  "description": "Complete taxonomy for SQL generation",
  "taxonomy": {
    "name": "SQL Operations",
    "children": [
      {
        "name": "Selection",
        "complexity_base": 1,
        "children": [
          {"name": "Simple WHERE", "complexity_base": 1},
          {"name": "Compound WHERE", "complexity_base": 2}
        ]
      }
    ]
  },
  "expected_node_count": 25,
  "expected_max_depth": 4
}
\end{lstlisting}

\subsection{Training example fixtures}

\begin{table}[htbp]
\centering
\caption{Training example fixture requirements.}
\label{tab:training-fixtures}
\footnotesize
\begin{tabularx}{\linewidth}{@{}l l X@{}}
\toprule
\textbf{Category} & \textbf{Count} & \textbf{Coverage} \\ \midrule
Easy examples & 20 & Single-table, simple predicates \\
Medium examples & 20 & Joins, grouping, ordering \\
Hard examples & 10 & CTEs, window functions, subqueries \\
Invalid examples & 10 & Should be rejected by critic \\
Edge cases & 10 & Unusual HANA types, large schemas \\
\bottomrule
\end{tabularx}
\end{table}

\section{Common fixtures}

\subsection{Mock LLM responses}

For deterministic testing without live LLM:

\begin{lstlisting}[basicstyle=\ttfamily\small,caption={Mock LLM response fixture.}]
{
  "fixture_id": "LLM-MOCK-001",
  "prompt_pattern": "Extract.*invoice.*fields",
  "response": {
    "content": "{\"seller\": {...}, \"totals\": {...}}",
    "usage": {"prompt_tokens": 500, "completion_tokens": 200},
    "finish_reason": "stop"
  },
  "latency_ms": 1500
}
\end{lstlisting}

\subsection{Schema copies}

All JSON Schema files \textbf{must} be copied to \texttt{tests/fixtures/common/schemas/}
to enable offline validation testing. Update procedure:

\begin{enumerate}
  \item Run \texttt{scripts/sync-test-schemas.sh} after schema changes
  \item Verify all fixtures still validate against updated schemas
  \item Commit schema copies alongside schema source
\end{enumerate}

\section{Data sanitisation requirements}

All fixtures derived from production data \textbf{must} be sanitised:

\begin{table}[htbp]
\centering
\caption{Data sanitisation rules.}
\label{tab:sanitisation}
\footnotesize
\begin{tabularx}{\linewidth}{@{}l X@{}}
\toprule
\textbf{Field Type} & \textbf{Sanitisation Method} \\ \midrule
Company names & Replace with ``ACME Corp'', ``Test Vendor'', etc. \\
Tax IDs & Generate valid-format fake numbers \\
Amounts & Randomise within same order of magnitude \\
Dates & Shift by random offset (preserve day-of-week) \\
Account numbers & Replace with sequential test numbers \\
Personal names & Replace with ``John Doe'', ``Jane Smith'', etc. \\
Addresses & Replace with generic test addresses \\
\bottomrule
\end{tabularx}
\end{table}

\section{Fixture validation}

\subsection{Schema validation}

All fixtures \textbf{must} pass schema validation:

\begin{lstlisting}[basicstyle=\ttfamily\small]
pytest tests/fixtures/validate_fixtures.py -v
\end{lstlisting}

\subsection{Coverage verification}

Fixture coverage is measured by:
\begin{itemize}
  \item Schema field coverage: Every schema field exercised
  \item Enum value coverage: Every enum value used at least once
  \item Edge case coverage: Boundary values for all numeric fields
  \item Error path coverage: All error codes triggerable
\end{itemize}

\section{Implementation checklist}

\begin{itemize}
  \item[\texttt{[ ]}] Arabic invoice fixtures (55 total per Table~\ref{tab:invoice-fixtures})
  \item[\texttt{[ ]}] Arabic OCR fixtures (all confidence levels)
  \item[\texttt{[ ]}] TB extract fixtures (33 total per Table~\ref{tab:tb-fixtures})
  \item[\texttt{[ ]}] TB variance fixtures (all anomaly types)
  \item[\texttt{[ ]}] Simula taxonomy fixtures (complete tree structure)
  \item[\texttt{[ ]}] Simula training example fixtures (70 total)
  \item[\texttt{[ ]}] Mock LLM response fixtures (all major prompt types)
  \item[\texttt{[ ]}] Schema copies in fixtures directory
  \item[\texttt{[ ]}] All fixtures pass schema validation
  \item[\texttt{[ ]}] Data sanitisation verified (no PII)
\end{itemize}
